\item Dos barras metálicas se fabrican de invar y una tercera barra se elabora de aluminio. A 0 °C, cada una de las tres barras se taladra con dos orificios separados por 40.0 cm. A través de los orificios se ponen clavijas para ensamblar las barras en un triángulo equilátero, como en la figura P6.31. 
(a) Primero ignore la expansión del invar. Encuentre el ángulo entre las barras de invar como función de la temperatura Celsius. 
(b) ¿Su respuesta es precisa tanto para temperaturas negativas como positivas? 
(c) ¿Es precisa para 0 °C? 
(d) Resuelva el problema de nuevo e incluya la expansión del invar. El aluminio se funde a 660 °C y el invar a 1 427 °C. Suponga que los coeficientes de expansión tabulados son constantes. ¿Cuáles son (e) el mayor y (f) el menor ángulos obtenibles entre las barras de invar?
